\documentclass[11pt,a4paper]{article}
\usepackage[T1]{fontenc}
\usepackage[utf8]{inputenc}
\usepackage{amsmath}
\usepackage{amssymb}
\usepackage{graphicx}
\usepackage[margin=2.5cm]{geometry}
\usepackage[hidelinks]{hyperref}

% Standalone draft of the two closing chapters. The counters continue the
% numbering of the report after the performance chapter, so that section
% numbers match the ones the chapters will carry once inserted.
\setcounter{section}{5}
\setcounter{table}{10}
\setcounter{figure}{31}
\setcounter{equation}{15}

\begin{document}

\section{Conclusions and future work}

This project developed a finite element framework for partial differential
equations posed on metric graphs and applied it to a problem of current
interest: the Fisher--Kolmogorov model of neurodegeneration on the brain
connectome. The library is built around a compact core consisting of a graph
topology with metric connections, a piecewise linear discretisation along
each connection, and a common workflow for coefficient evaluation, assembly,
solution and export. The benchmarks of Section 3 validated this framework on
problems with known solutions: stationary and time-dependent problems on the
star graph, showing the expected first-order accuracy of Backward Euler and
second-order accuracy of Crank--Nicolson in time, eigenvalue problems on the
star, graphene-like and binary tree graphs against the reference spectra of [6],
and the comparison between the metric-graph operator and the combinatorial
graph Laplacian.\\

The study of Section 4 used the library as a laboratory for the
Fisher--Kolmogorov equation on the 83-region connectome of [9]. Three findings
stand out. First, at one element per connection the metric-graph formulation
reproduces the behaviour of the nodal network models while remaining
computationally inexpensive, and refinement inside the connections remains
available whenever sharper fronts require it. Second, the balance between
reaction and transport, summarised by the Damk\"ohler number, governs both the
physical regime and the numerical robustness of the model, and the lumped
mass matrix is the appropriate choice in the reaction-dominated regime.
Third, connectivity alone orders the four lobes only partially: with the
regionally heterogeneous reaction rates of [8] transferred to our connectome,
the lobes cross the staging thresholds in the clinical order, suggesting a
division of roles between the transport pathways and the regional reaction
rates.\\

The performance analysis of Section 5 closed the loop between the mathematics
and the implementation. The unexpected cost of the sequential baseline was
traced to the fill-in generated by the default matrix ordering. The natural
numbering of the unknowns, with connection interiors first and graph vertices
last, exposes a bordered structure that can be exploited through static
condensation. By eliminating the interior unknowns locally, the global problem
is reduced to the 83 original graph vertices, while the remaining work becomes
naturally parallel over the connections. On a four-core laptop, combining
this reduction with parallel execution lowers the per-step cost of the finest
connectome discretisation by a factor of about thirty with respect to the best
full-system factorisation. Each intermediate optimisation was verified against
the sequential solver up to round-off error.\\

Several directions remain open.
\begin{itemize}
\item The MPI version was measured only on a single shared-memory node, where
the OpenMP implementation is preferable. Since the main purpose of MPI is to
enable distributed-memory execution, a multi-node scalability study is the
natural next step.

\item Static condensation currently accelerates the semi-implicit scheme.
Extending it to the fully implicit scheme means applying the same reduction
inside every Newton iteration, whose matrices share the sparsity structure
of the semi-implicit one.

\item A per-step cost of about two milliseconds at 143593 unknowns makes
many-run studies practical; parameter calibration and the uncertainty
quantification of [8] are the settings in which the speedup would be most
valuable.

\item The data pipeline already reconstructs the fine graph with 1015 nodes
from which the 83 regions are aggregated; repeating the study at the finer
parcellation would test the robustness of the staging results to the
aggregation.

\item Richer disease models from the literature, with several protein species
or patient-specific parameters, can be hosted on the same discretisation layer
without changes to its structure.
\end{itemize}

The quantitative results reported throughout this work are backed by stored
and regenerable benchmark records, as described in the next section.
Together with the validation of each optimisation against the reference
solver, this reproducibility is an important outcome of the project itself.\\

\section{C++ implementation}

\subsection{Repository reference}

The project is delivered as a single repository whose main entries separate
the library, the experiments and their records:

\begin{quote}
\begin{footnotesize}
\begin{verbatim}
FEMG/            the library: headers and sources of the problem classes
test_problem/    one driver per experiment family (parabolic, spectral,
                 fisher_kolmogorov, performance)
benchmarks/      one numbered record per experiment of the report, each
                 with a README, the commands used and the stored results
scripts/         benchmark runners, figure generation and verification
data/            graph files and the connectome preparation pipeline
verification/    internal correctness checks of the solver
report/          this document
\end{verbatim}
\end{footnotesize}
\end{quote}

\noindent
The records are the organising principle: every figure and every number of
the report is produced from files stored under \texttt{benchmarks/}, and
each record documents how to regenerate them.

\subsection{Dependencies}

The library is written in C++17 and built with CMake (3.16 or later). It
depends on Eigen 3.4 for sparse and dense linear algebra and on the Boost
Graph Library (1.76, the version provided by the mk toolchain of the
course) for the graph topology. OpenMP and MPI are optional: the
performance executables of Section 5 are added to the build only when CMake
finds them and the library itself requires neither. Post-processing uses
Python 3 with NumPy, Matplotlib and pandas. The default build type is
Release with \texttt{-O2}.

\subsection{Computational approach to differential problems}

All problems share the abstract base class
\texttt{quantum\_graph\_problem}. The topology is a Boost undirected
adjacency list whose vertex property carries the coordinates used for
visualisation and whose edge property carries the metric data of the
connection, its length and its number of finite element cells.
\texttt{init()} reads the graph file and builds the numbering of the
unknowns, interior nodes of each connection first and the original vertices
last, the layout on which Sections 5.2 and 5.3 rely. The base class stores
the shared discrete objects, the discrete Hamiltonian $H$, the mass matrix
$M$ and the solution vector, and prescribes the workflow as four virtual
steps: \texttt{set\_coefficients()}, \texttt{assembly()}, \texttt{solve()}
and \texttt{sol\_export()}.

Three concrete classes derive from it. \texttt{parabolic\_problem}
implements the $\theta$-method for the heat equation with Dirichlet
conditions on selected vertices and source and initial data supplied as
callables; as stated in Section 3.1, the elliptic benchmarks are stationary
solutions computed through this parabolic pipeline.
\texttt{spectral\_problem} solves the generalised eigenvalue problem
$Hx=\lambda Mx$ for a requested number of modes and also exposes the
eigenvalues of the combinatorial Laplacian for the comparison of Section
3.4.3. \texttt{fisher\_kolmogorov\_problem} implements the two time schemes
of Table 5, Backward Euler with a damped and safeguarded Newton iteration
and the semi-implicit scheme of [8], with one diffusion coefficient per
connection, one reaction rate per region and optional mass lumping. The
condensed solver of Section 5 is deliberately kept outside the library, as
a header-only engine under \texttt{test\_problem/performance/} that replays
the semi-implicit step without modifying the validated classes.

\subsection{Input data}

Graphs are plain text files. The header carries the number of vertices and
connections; an optional block with two coordinates per vertex follows, and
each connection is described by its two endpoint indices, its length and
its number of cells. The same reader accepts both the coordinate-free and
the coordinate-carrying variant. The \texttt{data/} directory provides the
analytical geometries of Section 3, the interval, the four-pointed star,
the graphene-like graph and the binary tree, at several resolutions. The
connectome enters through a documented preparation pipeline: the public
files of the Budapest Reference Connectome [10] are reduced to the
83-region connectome of [9] by the scripts under \texttt{scripts/}, which
also emit the per-connection weights and the refined meshes, from one to
128 elements per connection, used in Sections 4 and 5. Every derived file
is regenerable from the raw download.

\subsection{Setting coefficients}

Physical parameters travel through constructors and setters. The parabolic
constructor takes the diffusion coefficient, the final time, the time step
and $\theta$; sources, initial data and vertex Dirichlet data are function
objects evaluated on the metric coordinate of each connection. The
Fisher--Kolmogorov class accepts one diffusion coefficient per connection,
in Section 4 the normalised fibre weights, and one reaction rate and one
initial value per region, which it interpolates linearly along the
connections; further setters select the time scheme, the mass lumping, the
Newton parameters and the output stride. After the
setters, \texttt{set\_coefficients()} validates and stores everything the
assembly needs.

\subsection{Assembly procedures}

\texttt{assembly()} loops over the connections and, on each connection,
over its cells. The elemental blocks of the piecewise linear basis on a
cell of size $h$ are accumulated as triplets and assembled into compressed
Eigen sparse matrices: the consistent mass block, its lumped variant when
mass lumping is selected, and the stiffness block scaled by the diffusion
coefficient of the connection. The nonlinear reaction of the
Fisher--Kolmogorov model is integrated by Gauss quadrature applied to the
nodal interpolants of the rate and of the state, with a vertex rule
consistent with the lumped mass when lumping is on. The Kirchhoff
conditions at the vertices need no explicit treatment, since they arise
naturally from the weak form when the test functions are continuous across
a vertex; Dirichlet data are imposed algebraically on the selected
vertices.

\subsection{Solvers}

The linear heat equation factorises its $\theta$-method matrix once, before
the time loop, with the sparse LU of Eigen and reuses the factorisation at
every step. The Fisher--Kolmogorov schemes refactorise: the semi-implicit
scheme at every step, because the extrapolated reaction changes the matrix,
and the fully implicit scheme at every iteration of the damped and
safeguarded Newton method described in Section 4. The eigenvalue problems
of Section 3 are small enough for a dense generalised symmetric
eigensolver. The performance study of Section 5 adds the measured
alternatives: the symmetric $LDL^{T}$ factorisation on the natural
ordering, the per-connection static condensation with its dense
$83\times 83$ interface and the OpenMP, MPI and hybrid variants of the
condensed step.

\subsection{Exporting results}

\texttt{sol\_export()} writes ParaView files. Time-dependent problems
produce a \texttt{solution.pvd} collection referencing one
\texttt{solution\_XXXX.vtp} polyline file per stored step, with a
configurable stride to bound the output volume; the spectral problems write
the domain and one file per eigenfunction, together with a CSV of the
eigenvalues. The drivers add their own CSV summaries, timings, staging
times and regional means, which are the inputs of the figures.

\subsection{Result postprocessing}

Post-processing is Python. Each figure of the report is produced by one
script under \texttt{scripts/}, reading stored outputs, in almost all cases
the results of the corresponding benchmark record; shared modules fix the
visual conventions, so the figures of every chapter share one style. The
records close the loop. \texttt{scripts/run-benchmark.sh} rebuilds the
executable of one benchmark, reruns it and refreshes its stored results,
and each record's README documents the reading of its experiment;
\texttt{scripts/verify-figures.py} re-derives from the stored outputs
several hundred of the quantities quoted in the report, from mesh sizes and
factor fills to staging times and validation differences, and fails when
any of them drifts. The figures of Section 5 are generated in the same way from
the CSV records of the performance studies.

\subsection{Example of code usage}

The driver of Test E3 (Section 3.1.3) shows the complete arc of the
interface: construction with the physical and time parameters, data
supplied as function objects, initialisation from a graph file, coefficient
evaluation, assembly, solution and the error check against the exact
solution.

\begin{quote}
\begin{footnotesize}
\begin{verbatim}
#include "parabolic_problem.hpp"
#include "star_sine_functions.hpp"
#include <iostream>
#include <memory>
#include <string>

int main(int argc, char *argv[]) {
    const double diffusion = 1.0;
    const double T = 1.0;
    const double deltat = 0.005;
    const double theta = 1.0;

    std::string graph_file = "data/star_4.txt";
    char *local_argv[] = {argv[0], graph_file.data()};

    femg::parabolic_problem problem(diffusion, T, deltat, theta);
    problem.set_initial_condition(std::make_shared<StarSineInitialCondition>());
    problem.set_source_function(std::make_shared<StarSineForcingTerm>());
    problem.set_dirichlet_vertices(
        {0, 1, 2, 3, 4},
        [](std::size_t, double) {
            return 0.0;
        });
    problem.set_output_directory("output/visualization/star_sine");

    if (argc >= 2) {
        problem.init(argc, argv);
    } else {
        problem.init(2, local_argv);
    }

    problem.set_coefficients();
    problem.assemble_matrices();
    problem.print_matrix_summary(std::cout);
    problem.solve();

    const StarSineExactSolution exact_solution;
    const double error_l2 = problem.compute_l2_error(exact_solution);
    std::cout << "Final L2 error against u = sin(2*pi*x): "
              << error_l2 << "\n";

    return 0;
}
\end{verbatim}
\end{footnotesize}
\end{quote}

\noindent
Building and running the example requires three commands:

\begin{quote}
\begin{footnotesize}
\begin{verbatim}
cmake -S . -B build
cmake --build build
./build/test_heat_star_sine
\end{verbatim}
\end{footnotesize}
\end{quote}

\noindent
The solution lands in \texttt{output/visualization/star\_sine/solution.pvd},
ready for ParaView. The same pattern, a constructor, the setters,
\texttt{init}, coefficients, assembly, solve and export, is the whole
interface a new experiment needs.

\end{document}
